\documentclass[11pt]{article}
\usepackage[a4paper,margin=1in]{geometry}
\usepackage{fontspec}
\usepackage[boldfont=false]{xeCJK}
\setCJKmainfont{FandolSong-Regular.otf}
\usepackage{amsmath}
\usepackage{amssymb}
\usepackage{array}
\usepackage{booktabs}
\usepackage{graphicx}
\usepackage{adjustbox}
\usepackage{enumitem}
\setlist[itemize]{topsep=2pt,partopsep=0pt,parsep=0pt,itemsep=2pt,leftmargin=1.4em}
\setlist[enumerate]{topsep=2pt,partopsep=0pt,parsep=0pt,itemsep=2pt,leftmargin=1.6em}
\usepackage{parskip}
\usepackage{fvextra}
\fvset{breaklines=true,breakanywhere=true,fontsize=\small}
\setlength{\parindent}{0pt}

\begin{document}

\begin{center}
  {\LARGE\bfseries Biological molecules}\\[0.35em]
  {\large A Level Biology}
\end{center}
\vspace{0.6em}

\section*{Testing for biological molecules}

Living things are built from four main kinds of large \textbf{molecule} 分子: \textbf{carbohydrates} 碳水化合物, \textbf{lipids} 脂质, \textbf{proteins} 蛋白质 and \textbf{nucleic acids} 核酸. You can use simple chemical tests to find out which kinds are present in a sample.

\begin{center}
\adjustbox{max width=\linewidth}{%
\begin{tabular}{llll}
\toprule
\textbf{Test} & \textbf{What it finds} & \textbf{Method} & \textbf{Positive result} \\
\midrule
Benedict's test & \textbf{reducing sugar} 还原糖 & add Benedict's solution and heat in a water bath & blue changes to green, yellow, orange, then brick-red \textbf{precipitate} 沉淀 \\
iodine test & \textbf{starch} 淀粉 & add orange-brown \textbf{iodine} 碘 solution & colour changes to blue-black \\
emulsion test & \textbf{lipid} & mix sample with ethanol, then pour into water & a white, cloudy \textbf{emulsion} 乳浊液 forms \\
\textbf{biuret} 双缩脲 test & \textbf{protein} & add biuret solution at room temperature & blue changes to purple \\
\bottomrule
\end{tabular}}
\end{center}

\subsection*{Semi-quantitative Benedict's test}

The normal Benedict's test only tells you "yes or no". A \textbf{semi-quantitative} 半定量 test gives a rough amount. First you \textbf{standardise} 标准化 the test: you run it on solutions of known \textbf{concentration} 浓度 and record the result for each. Then you can estimate an unknown by either:

\begin{itemize}
\item the \textbf{time to the first colour change} (more sugar changes colour faster), or

\item comparing the final colour to your set of colour standards.

\end{itemize}

\subsection*{Testing for non-reducing sugars}

Some sugars, such as sucrose, are \textbf{non-reducing sugar} 非还原糖: they give a negative Benedict's test. To detect them:

\begin{enumerate}
\item Do a normal Benedict's test first. It stays blue (no reducing sugar).

\item Take a fresh sample and add dilute \textbf{hydrochloric acid} 盐酸, then heat. This \textbf{acid hydrolysis} breaks the sugar into smaller reducing sugars.

\item Cool, then \textbf{neutralise} 中和 the acid with sodium hydrogencarbonate.

\item Now do the Benedict's test again. A brick-red colour shows a non-reducing sugar was present.

\end{enumerate}

\section*{Carbohydrates}

\subsection*{Monomers, polymers and macromolecules}

\begin{itemize}
\item a \textbf{monomer} 单体 is a small molecule that is a single unit.

\item a \textbf{polymer} 聚合物 is a long molecule made of many monomers joined together.

\item a \textbf{macromolecule} 大分子 is any very large molecule.

\end{itemize}

Sugars come in three sizes:

\begin{itemize}
\item a \textbf{monosaccharide} 单糖 is a single sugar unit, such as \textbf{glucose} 葡萄糖 and fructose.

\item a \textbf{disaccharide} 二糖 is two units joined, such as maltose and sucrose.

\item a \textbf{polysaccharide} 多糖 is many units joined into a polymer.

\end{itemize}

Glucose, \textbf{fructose} 果糖 and \textbf{maltose} 麦芽糖 are reducing sugars. \textbf{Sucrose} 蔗糖 is a non-reducing sugar.

\subsection*{Two ring forms of glucose}

Glucose has six carbon atoms and forms a ring. There are two ring forms. They differ only at carbon 1:

\begin{itemize}
\item in \textbf{α-glucose}, the –OH group on carbon 1 points \textbf{down}, below the ring.

\item in \textbf{β-glucose}, the –OH group on carbon 1 points \textbf{up}, above the ring.

\end{itemize}

This small difference decides which polysaccharide the glucose can build.

\subsection*{Joining and breaking sugars}

Monomers are joined by strong \textbf{covalent bonds} 共价键. When two sugars join, a \textbf{glycosidic bond} 糖苷键 forms between them. This happens by \textbf{condensation} 缩合: a molecule of water is removed each time a bond forms.

The reverse is \textbf{hydrolysis} 水解: a water molecule is added to break a glycosidic bond. This is why the non-reducing sugar test needs acid and heat — they hydrolyse sucrose into glucose and fructose.

\subsection*{Storage polysaccharides: starch and glycogen}

\textbf{Starch} is the energy 能量 store in plants. It is made of two polymers of α-glucose:

\begin{itemize}
\item \textbf{amylose} 直链淀粉 — a long, unbranched chain that coils into a spiral.

\item \textbf{amylopectin} 支链淀粉 — a chain with many side branches.

\end{itemize}

\textbf{Glycogen} 糖原 is the energy store in animals. It is like amylopectin but has even more branches, so it can be broken down quickly when energy is needed.

These stores suit their job well: they are compact, they are \textbf{insoluble} 不溶 (so they do not leave the cell), and they do not change the \textbf{water potential} 水势 of the cell (so they do not pull water in by \textbf{osmosis} 渗透). The many branches give many ends, so glucose can be added or removed fast.

\subsection*{Cellulose}

\textbf{Cellulose} 纤维素 is made of β-glucose. Because of the β form, every other glucose is flipped over, so the chains are long and straight. Many straight chains lie side by side and are held together by \textbf{hydrogen bonds} 氢键 into strong bundles called \textbf{microfibrils} 微纤丝. These give the plant \textbf{cell wall} 细胞壁 its strength and stop the cell bursting.

\section*{Lipids}

\subsection*{Triglycerides}

A \textbf{triglyceride} 甘油三酯 is the main fat or oil. It is \textbf{non-polar} 非极性 and \textbf{hydrophobic} 疏水 (it does not mix with water). It is made from one \textbf{glycerol} 甘油 molecule joined to three \textbf{fatty acids} 脂肪酸 by \textbf{ester bonds} 酯键. Each ester bond forms by condensation, so three water molecules are removed.

Fatty acids are of two kinds:

\begin{itemize}
\item \textbf{saturated} 饱和 — no carbon–carbon \textbf{double bonds} 双键; these fats are usually solid.

\item \textbf{unsaturated} 不饱和 — one or more double bonds; these oils are usually liquid.

\end{itemize}

Triglycerides make a good long-term energy store: they release about twice as much energy per gram as carbohydrates, they are insoluble, and they store little extra mass because they hold no water. Under the skin they also give \textbf{insulation} 隔热 and protect the organs.

\subsection*{Phospholipids}

A \textbf{phospholipid} 磷脂 is like a triglyceride, but one fatty acid is replaced by a \textbf{phosphate} 磷酸 group. This gives the molecule two ends with different behaviour:

\begin{itemize}
\item a \textbf{hydrophilic} 亲水 ("water-loving") \textbf{polar} 极性 phosphate head.

\item two hydrophobic ("water-fearing") fatty acid tails.

\end{itemize}

This split personality is why phospholipids form the membranes around cells.

\section*{Proteins}

\subsection*{Amino acids and the peptide bond}

Proteins are polymers of \textbf{amino acids} 氨基酸. Every amino acid has the same general structure around a central carbon atom: an \textbf{amino group} 氨基 (–NH₂), a \textbf{carboxyl group} 羧基 (–COOH), a hydrogen atom, and a variable \textbf{side chain} 侧链 (the R group). The R group is different in each amino acid.

Two amino acids join by condensation. The bond formed between the amino group of one and the carboxyl group of the next is a \textbf{peptide bond} 肽键, and a water molecule is removed. Many amino acids joined this way make a \textbf{polypeptide} 多肽. Adding water (hydrolysis) breaks a peptide bond.

\subsection*{Four levels of protein structure}

\begin{center}
\adjustbox{max width=\linewidth}{%
\begin{tabular}{ll}
\toprule
\textbf{Level} & \textbf{What it means} \\
\midrule
\textbf{primary structure} 一级结构 & the order of amino acids in the chain \\
\textbf{secondary structure} 二级结构 & local shapes — the \textbf{α-helix} 螺旋 and the \textbf{β-pleated sheet} 折叠片 — held by hydrogen bonds \\
\textbf{tertiary structure} 三级结构 & the whole chain folded into a precise 3-D shape \\
\textbf{quaternary structure} 四级结构 & two or more polypeptide chains joined into one protein \\
\bottomrule
\end{tabular}}
\end{center}

The folded shape is held together by four kinds of interaction between R groups:

\begin{itemize}
\item \textbf{hydrophobic interactions} 疏水作用 (non-polar R groups cluster away from water).

\item hydrogen bonding.

\item \textbf{ionic bonds} 离子键 (between charged R groups).

\item covalent bonding, including strong \textbf{disulfide bonds} 二硫键.

\end{itemize}

\subsection*{Globular and fibrous proteins}

\begin{itemize}
\item \textbf{globular proteins} 球状蛋白质 fold into a rounded shape, are usually \textbf{soluble} 可溶, and do jobs in the body (for example enzymes and haemoglobin).

\item \textbf{fibrous proteins} 纤维状蛋白质 form long strands, are usually insoluble, and give structure and support (for example collagen).

\end{itemize}

\subsection*{Haemoglobin — a globular protein}

\textbf{Haemoglobin} 血红蛋白 carries \textbf{oxygen} 氧气 in red blood cells. It has a quaternary structure made of four polypeptide chains: two alpha (α-globin) chains and two beta (β-globin) chains. Each chain holds a \textbf{haem group} 血红素. At the centre of each haem group is an \textbf{iron} 铁 atom, and this is where one oxygen molecule binds. Four chains mean one haemoglobin molecule can carry four oxygen molecules.

\subsection*{Collagen — a fibrous protein}

\textbf{Collagen} 胶原蛋白 gives strength to skin, \textbf{tendons} 肌腱, bone and blood vessel walls. One collagen molecule is three polypeptide chains wound tightly around each other in a triple strand, held by hydrogen bonds. Many of these molecules lie side by side, slightly staggered, and are cross-linked into thick \textbf{fibres} 纤维. The staggered, cross-linked arrangement makes collagen very strong when pulled.

\section*{Water}

Water is a small molecule, but its two O–H bonds are \textbf{polar}: the oxygen end is slightly negative and the hydrogen ends are slightly positive. So one water molecule attracts its neighbours, forming weak hydrogen bonds between them. These hydrogen bonds explain water's useful properties:

\begin{itemize}
\item \textbf{solvent action} — water is a good \textbf{solvent} 溶剂, so many substances dissolve in it. This lets reactions happen and lets substances be carried around the body.

\item \textbf{high specific heat capacity} 比热容 — water needs a lot of energy to warm up, so its temperature stays steady. This protects living things from quick temperature changes.

\item \textbf{latent heat of vaporisation} 汽化潜热 — water needs a lot of energy to \textbf{evaporate} 蒸发. So when water evaporates (for example as sweat dries), it carries away a lot of heat and cools the body.

\end{itemize}


\end{document}
